\subsection{Regular-prime freeness and prescribed support matching}

\begin{result}{P.15}{Proved}
{The available embedding theorems retain a separate local matching hypothesis}
The regular-prime vanishing in E.209 does not by itself verify the
hypotheses of the restricted-ramification embedding theorems below.
In particular, for
\[
 A=\mathbb F_5(3),\qquad S=\{5,\infty\},
 \qquad T=\{11,31\},
\]
Liu's finite obstruction group satisfies
\[
       B^{S\cup T}_{S\setminus T}(\mathbb Q,A)\ne0,
       \qquad
       \dim_{\mathbb F_5}B^{S\cup T}_{S\setminus T}(\mathbb Q,A)
          \ge1.                                                \tag{P.15.1}
\]
This holds although all the fixed-$S$ second-cohomology groups
in E.209 vanish at five. The calculation is a distinction between
two hypotheses. It does not assert that a discrepancy tuple coming
from the ordinary point sections in B.26 is nonzero in this group.

The regular-prime projectivity construction of P.4, Schmidt's
free-product theorem, and the solvable embedding theorem of
Jarden--Ramiharimanana do not supply the missing simultaneous
point matches for one ordinary pair with infinite support.
The precise comparison maps and the retained restrictions are
given below. No counterexample to A.2 or A.3 is asserted.
\end{result}

\begin{proof}[Theorem hypotheses and comparison]
\emph{The projective source and the actual ramified source.}
Let $\Omega$ be the maximal pro-five extension of
$K=\mathbb Q(\zeta_5)$ unramified outside five, and put
$\Gamma=\operatorname{Gal}(\Omega/\mathbb Q)$.
Let $P$ be the actual free pro-five geometric fundamental group
of the tripod, based at the rational unit tangent at zero.
The subgroup $\operatorname{Gal}(\Omega/K)$ is free pro-five
by the regular-prime theorem used in P.4
\citep[Lemma 3.2, author p.~14]{LitHoshi2010}.
Its index in $\Gamma$ is four. Thus $\operatorname{cd}_5\Gamma\le1$,
which solves the pro-five-kernel problem
\[
 P\rtimes\Gamma\longrightarrow P^{\rm ab}\rtimes\Gamma
\]
for an abelian cocycle defined on $\Gamma$.
The pointed geometric action does factor through $\Gamma$.
An ordinary completed abelian pair with nonzero valuation at a
prime $q\ne5$ does not: its tame inertia value in at least one
coordinate is the nonzero ordinary integer $e_q\in\mathbb Z_5$.
The source map $G_{\mathbb Q}\to\Gamma$ kills that inertia.
Even if $e_q$ is divisible by a large power of five, it is not
zero in $\mathbb Z_5$. Therefore the graph to be lifted is a
map on $G_{\mathbb Q}$, and cannot be replaced by the graph of
a cocycle on this projective quotient.

\emph{A free product inside a relative kernel.}
For this paragraph use $S_1,T_1$ for finite sets of places of a
global field $k$, and let $M$ be disjoint from them and have
Dirichlet density zero. For odd $p\ne\operatorname{char}(k)$,
Schmidt's theorem produces a finite auxiliary set $S_0$ disjoint
from $S_1\cup T_1\cup M$. With
\[
 L=k^{T_1}_{S_1\cup S_0}(p),\qquad
 L_0=k^{S_1\cup T_1}_{S_0}(p),
\]
its exact free-product assertion, in the category of pro-$p$ groups, is
\[
 \mathop{\ast}_{v\in S_1(L_0)}G_{k_v}(p)
       \xrightarrow{\ \sim\ }\operatorname{Gal}(L/L_0).
                                                               \tag{P.15.2}
\]
Here superscripts prescribe complete splitting and subscripts
prescribe the allowed ramification. Local maximal pro-$p$
extensions at $S_1\cup S_0$ are realized
\citep[Theorem 1.1(ii)--(iii), author p.~2;
Proposition 7.1 and its proof, pp.~19--20]{RootSchmidt2010}.

The group on the right of P.15.2 is the kernel in
\[
 1\longrightarrow\operatorname{Gal}(L/L_0)
 \longrightarrow\operatorname{Gal}(L/k)
 \longrightarrow\operatorname{Gal}(L_0/k)\longrightarrow1.
                                                               \tag{P.15.3}
\]
It is not asserted to be a free factor of the middle group.
Prescribing homomorphisms on the local factors determines a
homomorphism on this kernel only when the associated family is
continuous. Extension to the middle group must also respect
its conjugation action: for $n$ in the kernel and $g$ in the
middle group one needs
\[
                  \phi(gng^{-1})=\phi(g)\phi(n)\phi(g)^{-1}.
                                                               \tag{P.15.4}
\]
The old Frobenius images and the actual arithmetic action enter
this equation. P.15.2 does not prescribe them, and it does not
specify peripheral images for inertia at $S_0$. Reapplying the
theorem to larger finite sets changes the extension in P.15.3;
the theorem does not supply compatible point-section maps on
these varying extensions. In particular, its ability to avoid a
density-zero set when choosing $S_0$ is not a theorem prescribing
all section values on that set.

\emph{Solvable embedding with finitely many prescribed solutions.}
Let $K_0\subset K_1$ be a finite Galois extension of global
fields. Jarden--Ramiharimanana consider a finite embedding problem
\[
 G_{K_0}\longrightarrow\operatorname{Gal}(K_1/K_0),\qquad
 1\longrightarrow H\longrightarrow E
 \longrightarrow\operatorname{Gal}(K_1/K_0)\longrightarrow1,
\]
where $H$ is solvable,
$\operatorname{char}(K_0)\nmid |H|$, and
\[
                     \gcd(|H|,|\mu(K_1)|)=1.                  \tag{P.15.5}
\]
Assuming local solvability everywhere, their Theorem A gives a
proper solution agreeing, up to $H$-conjugacy, with specified
solutions at a finite set containing the old ramification. It
allows at most $\Omega(|H|)$ additional ramified primes, where
$\Omega$ counts prime factors with multiplicity
\citep[Theorem A, author pp.~3--4]{LitJardenRamiharimanana2018}.

In the application retaining the residual cyclotomic character,
the finite quotient field $K_1$ contains $\mathbb Q(\zeta_5)$.
A nontrivial five-primary central layer has $5\mid |H|$, so
P.15.5 fails. Passing to $\mathbb Q(\zeta_5)$ as the base does
not remove this failure. Discarding that cyclotomic quotient
would change the arithmetic embedding problem. Even in a
different problem satisfying P.15.5, the theorem's conclusion
places no specified tripod peripheral generator at each of its
new ramification primes.

\emph{The restricted-ramification obstruction is a different group.}
Write $A^D=\operatorname{Hom}(A,\mu_5)$, and use a superscript
$\vee$ for Pontryagin duality. In the notation of
\citet[Definition 3.1 and Lemma 3.2, author pp.~8--9]{LitLiu2024},
for sets $V\subseteq U$ the group $B^U_V(k,A)$ is the cokernel of
\[
 \left(\prod_{v\in V}'H^1(k_v,A)\right)
 \times\left(\prod_{v\notin U}H^1_{\rm nr}(k_v,A)\right)
 \longrightarrow H^1(k,A^D)^\vee,                              \tag{P.15.6}
\]
using the sum of local Tate pairings. The prime on a product
means unramified at all but finitely many places. Lemma 3.7,
author pp.~11--12, assumes that the predecessor and the finite
simple coefficient module factor through $G_{k,S}^{T}$,
that the module's trivializing extension is solvable, and that
$B^{S\cup T}_{S\setminus T}(k,A)=0$. Under these hypotheses,
local solvability gives a solution unramified outside $S$ and
split at $T$. Its proof uses the exact correction sequence
\[
 \begin{gathered}
 H^1(k,A)\longrightarrow
 \bigoplus_{v\notin S\cup T}H^1(I_v,A)^{G_{k_v}}
       \oplus\prod_{v\in T}'H^1(k_v,A)\\
 \longrightarrow B^{S\cup T}_{S\setminus T}(k,A)
 \longrightarrow0.
 \end{gathered}
 \tag{P.15.7}
\]
Thus this vanishing is precisely an additional approximation
hypothesis. The statement does not identify $T$ with arbitrary
nontrivial local sections at the old support; its $T$ consists of
places where the predecessor is already split. Nor does the
theorem assert chosen local solutions at every place of $S$.

We now compute P.15.1 directly from P.15.6. Set
\[
 U=S\cup T=\{5,11,31,\infty\},\qquad V=S,
 \qquad A^D=\mathbb F_5(-2)\simeq\mathbb F_5(2).
\]
Both $11$ and $31$ are one modulo five. Hence $A$ is indeed
a module for $G_{\mathbb Q,S}^{T}$: its trivializing field is
$\mathbb Q(\zeta_5)$, which is unramified outside $S$ and
split at both primes of $T$.

The exact annihilator of unramified first cohomology at a
prime outside $U$ is unramified first cohomology with the dual
module \citep[I.2.6, pp.~30--31]{CohMilneADT}. All these primes
have residue characteristic different from five, and both
modules are unramified there. The annihilator of the full
local group at five is zero. The real term is zero because
five is odd. Dualizing P.15.6 consequently gives the exact
identification
\[
 \left(B^U_V(\mathbb Q,A)\right)^\vee
   =\ker\left(
       H^1(G_{\mathbb Q,U},\mathbb F_5(2))
        \longrightarrow H^1(\mathbb Q_5,\mathbb F_5(2))
         \right).                                             \tag{P.15.8}
\]
There is no extra local condition at $11$ or $31$ in this
kernel. At primes outside $U$ the conditions are unramified,
not locally trivial. This distinction fixes the $S/T$ convention.

The finite coefficient sequence and E.209--E.210 give
\[
 H^1(G_{\mathbb Q,S},\mathbb F_5(2))=
 H^2(G_{\mathbb Q,S},\mathbb F_5(2))=0.
\]
The finite purity sequence for adjoining $11$ and $31$ is
therefore an isomorphism on the displayed residue map:
\[
 H^1(G_{\mathbb Q,U},\mathbb F_5(2))
    \xrightarrow{\ \sim\ }
 H^0(\mathbb F_{11},\mathbb F_5(1))
       \oplus H^0(\mathbb F_{31},\mathbb F_5(1))
    =\mathbb F_5^2.                                           \tag{P.15.9}
\]
This is the same finite supports sequence used in E.202;
it makes no assertion about an inverse limit over newly
introduced ramification.
At five, $H^0(\mathbb Q_5,\mathbb F_5(2))=0$ since the
residual cyclotomic image has order four. Local duality gives
$H^2(\mathbb Q_5,\mathbb F_5(2))=0$, since its dual weight
is $-1$ and likewise has no invariants. The local Euler
characteristic then gives
\[
           \dim_{\mathbb F_5}H^1(\mathbb Q_5,\mathbb F_5(2))=1.
                                                               \tag{P.15.10}
\]
The kernel in P.15.8 has dimension at least one by
P.15.9--P.15.10. This proves P.15.1 without guessing a
local class, a Frobenius representative, or the image of an
ordinary point's discrepancy tuple.

Finally, at a fixed integral weight $w\ge2$, every difference
between two global lifts with the same literal predecessor
belongs to $H^1(G_{\mathbb Q},\mathbb Z_5(w))$.
At $q\ne5$ its tame inertia coordinate $v$ satisfies
$(q^w-q)v=0$, so $v=0$. Adding auxiliary ramification to a
finite coefficient problem can enlarge its correction space;
it does not enlarge this integral space. The same observation
does not prohibit changing the eventual ordinary pair during a
new construction. Such a construction must additionally arrange
the specified peripheral inertia and compatible local point
paths, which none of the three quoted theorems supplies.
The remaining condition is the actual localization image in
B.26, now computed at regular primes by E.210.
\end{proof}
